\chapter{Исследовательский раздел}

\section{Цель и методика исследования}

Целью раздела является оценка производительности аналитических запросов
при изменении объёма данных, конфигурации индексов, состояния кэша и уровня
параллельной нагрузки. Измерения выполняются для ключевых механизмов системы:
предвычисленной витрины \texttt{player\_season\_stats}, индексов B-tree,
Redis-кэша и асинхронного API-сервера.

Замеры автоматизированы Python-скриптом: SQL-запросы выполняются через
asyncpg, HTTP-запросы к API~--- через httpx. PostgreSQL~15, Redis~7
и API-сервер развёрнуты в Docker-контейнерах на одном хосте:
CPU~--- Intel Core i5-1235U (10~ядер, 3{,}3~ГГц), ОЗУ~--- 16~ГБ DDR4,
хранилище~--- NVMe SSD. Таблица фактов \texttt{game\_player\_stats}
содержит 159\,652~строк за пять сезонов NBA.

Для SQL-запросов выполнялось 10~прогревочных запусков и 10\,000~рабочих
повторов. По рабочим замерам рассчитывались медиана, среднее значение
и 95-й перцентиль. Выбросы отдельно не отсекались. Для проверки кэширования
измерялись Cache Miss (30~запросов со сбросом кэша перед каждым запросом)
и Cache Hit (10\,000~запросов после прогрева кэша). Нагрузочные замеры выполнены
для 1, 10, 50 и 100~параллельных запросов, по 10~раундов на каждый уровень.
Ошибкой считается HTTP-код не ниже 500, таймаут или отказ соединения.

\section{Влияние объёма данных на производительность}

Для оценки влияния объёма данных сопоставлены два способа получения рейтинга
лидеров: прямая агрегация по таблице фактов и выборка из предвычисленной
витрины. Объём изменялся накопительно: от одного сезона до пяти сезонов.
Эффект витрины оценивается кратностью ускорения~--- отношением медианного
времени прямой агрегации к медианному времени выборки из витрины на одном
и том же наборе сезонов.

Запрос прямой агрегации по таблице фактов приведён
в листинге~\ref{lst:research-heavy-agg}.

\begin{lstlisting}[language=SQL, label=lst:research-heavy-agg,
  caption=Прямой агрегирующий запрос к таблице фактов]
SELECT gps.player_id,
       SUM(gps.points)        AS total_pts,
       COUNT(*)               AS games,
       AVG(gps.points::float) AS avg_pts
FROM game_player_stats gps
JOIN games g ON g.game_id = gps.game_id
WHERE g.season_id = ANY($1::int[])
GROUP BY gps.player_id
ORDER BY total_pts DESC
LIMIT 20;
\end{lstlisting}

Запрос к предвычисленной витрине приведён
в листинге~\ref{lst:research-vitrina-query}.

\begin{lstlisting}[language=SQL, label=lst:research-vitrina-query,
  caption=Запрос к предвычисленной витрине]
SELECT player_id, season_id, per, avg_pts
FROM player_season_stats
WHERE season_id = ANY($1::int[])
ORDER BY per DESC NULLS LAST
LIMIT 20;
\end{lstlisting}

Результаты замеров приведены в таблице~\ref{tab:volume-fact-vitrina}.

\begin{table}[H]
\caption{Влияние объёма данных на прямую агрегацию и выборку из витрины}
\label{tab:volume-fact-vitrina}
\centering
\small
\begin{tabular}{|p{0.13\textwidth}|p{0.16\textwidth}|p{0.16\textwidth}
  |p{0.16\textwidth}|p{0.16\textwidth}|p{0.16\textwidth}|}
\hline
\textbf{Сезоны} & \textbf{Строк фактов} & \textbf{Строк витрины}
  & \textbf{Агрегация, медиана (мс)} & \textbf{Витрина, медиана (мс)}
  & \textbf{Ускорение, раз} \\
\hline
1       & 31\,311  & 586   & 18{,}58 & 0{,}44 & 42{,}2 \\
\hline
1--2    & 62\,853  & 1\,140 & 24{,}75 & 0{,}66 & 37{,}5 \\
\hline
1--3    & 95\,235  & 1\,721 & 25{,}25 & 0{,}78 & 32{,}4 \\
\hline
1--4    & 127\,607 & 2\,304 & 29{,}09 & 0{,}90 & 32{,}3 \\
\hline
1--5    & 159\,652 & 2\,897 & 35{,}58 & 1{,}15 & 30{,}9 \\
\hline
\end{tabular}
\end{table}

Зависимость медианного времени выполнения от объёма данных показана
на рисунке~\ref{fig:volume-fact-vitrina}.

\begin{figure}[H]
\centering
\begin{tikzpicture}
\begin{axis}[
    width=0.88\textwidth,
    height=6.5cm,
    xlabel={Количество строк фактов},
    ylabel={Медианное время выполнения (мс)},
    xmin=25000, xmax=165000,
    ymin=0, ymax=38,
    scaled x ticks=false,
    xtick={31311,62853,95235,127607,159652},
    xticklabels={31\,311,62\,853,95\,235,127\,607,159\,652},
    x tick label style={font=\scriptsize, rotate=25, anchor=east},
    ymajorgrids=true,
    grid style=dashed,
    legend style={
      at={(0.5,1.14)},
      anchor=south,
      legend columns=2,
      font=\small,
      draw=none
    },
    mark size=3pt,
]
\addplot+[mark=*, thick, color=black] coordinates {
    (31311, 18.58)
    (62853, 24.75)
    (95235, 25.25)
    (127607, 29.09)
    (159652, 35.58)
};
\addplot+[mark=square*, thick, color=black, dashed] coordinates {
    (31311, 0.44)
    (62853, 0.66)
    (95235, 0.78)
    (127607, 0.90)
    (159652, 1.15)
};
\legend{Прямая агрегация, Витрина}
\end{axis}
\end{tikzpicture}
\caption{Влияние объёма данных на время аналитической выборки}
\label{fig:volume-fact-vitrina}
\end{figure}

При росте таблицы фактов с 31\,311 до 159\,652~строк медиана прямой
агрегации увеличилась с 18{,}58 до 35{,}58~мс. Запрос к витрине выполняется
за 0{,}44--1{,}15~мс, поскольку использует готовые сезонные агрегаты.
На полном объёме данных витрина ускоряет выборку примерно в 31~раз.

Кратность ускорения максимальна на одном сезоне (42{,}2~раза) и плавно
снижается к пяти сезонам. Причина~--- в устройстве кумулятивного запроса
к витрине: одиночный сезон обслуживается одним диапазонным сканом индекса,
тогда как \texttt{season\_id = ANY(...)} сканирует несколько диапазонов
и пересортировывает строки по \texttt{per} через все сезоны. Из-за этих
накладных расходов время витрины с ростом числа сезонов увеличивается
пропорционально быстрее времени агрегации (в абсолютных величинах витрина
остаётся быстрее в десятки раз), поэтому кратность ускорения убывает,
оставаясь в пределах 31--42~раз.

\section{Влияние индексирования на селективную выборку}

Влияние индексов оценивалось на селективной выборке~--- получении полного
игрового лога одного игрока (листинг~\ref{lst:research-player-log}). Запрос
фильтрует таблицу фактов по \texttt{player\_id} и возвращает 410~строк
из 159\,652, то есть 0{,}26\,\% объёма. Индекс \texttt{idx\_gps\_player\_id}
спроектирован в подразделе~\ref{subsec:indexes}; его DDL-объявление приведено
в листинге~\ref{lst:indexes} и Приложении~А.

\begin{lstlisting}[language=SQL, label=lst:research-player-log,
  caption=Селективная выборка игрового лога игрока]
SELECT gps.game_id, gps.points, gps.assists,
       gps.rebounds_off, gps.rebounds_def
FROM game_player_stats gps
WHERE gps.player_id = $1
ORDER BY gps.game_id;
\end{lstlisting}

Сравнивались две конфигурации: выполнение без индекса по фильтруемому
столбцу (последовательный просмотр таблицы) и с индексом
\texttt{idx\_gps\_player\_id}. Результаты приведены
в таблице~\ref{tab:fact-indexes}.

\begin{table}[H]
\caption{Влияние индекса по столбцу фильтрации на селективную выборку}
\label{tab:fact-indexes}
\centering
\small
\begin{tabular}{|p{0.40\textwidth}|p{0.14\textwidth}|p{0.14\textwidth}|p{0.14\textwidth}|}
\hline
\textbf{Конфигурация} & \textbf{Медиана (мс)}
  & \textbf{Среднее (мс)} & \textbf{P95 (мс)} \\
\hline
Без индекса (Seq Scan) & 13{,}05 & 17{,}99 & 48{,}74 \\
\hline
\texttt{idx\_gps\_player\_id} (Bitmap Index Scan) & 2{,}46 & 3{,}44 & 11{,}08 \\
\hline
\end{tabular}
\end{table}

Сравнение времени выполнения без индекса и с индексом показано
на рисунке~\ref{fig:fact-indexes}.

\begin{figure}[H]
\centering
\begin{tikzpicture}
\begin{axis}[
    ybar,
    bar width=26pt,
    width=0.88\textwidth,
    height=6.6cm,
    ylabel={Время выполнения (мс)},
    symbolic x coords={Без индекса, С индексом},
    xtick=data,
    x tick label style={font=\small},
    enlarge x limits=1.2,
    ymin=0, ymax=56,
    ytick={0,10,20,30,40,50},
    ymajorgrids=true,
    grid style=dashed,
    nodes near coords,
    nodes near coords style={font=\scriptsize,
      /pgf/number format/fixed,
      /pgf/number format/precision=2},
    every node near coord/.append style={anchor=south, yshift=1pt},
    legend style={at={(0.5,1.10)}, anchor=south, legend columns=2,
      font=\small, draw=none},
]
\addplot+[fill=white, draw=black,
          postaction={pattern=north east lines}] coordinates {
    (Без индекса, 13.05)
    (С индексом, 2.46)
};
\addplot+[fill=black!15, draw=black] coordinates {
    (Без индекса, 48.74)
    (С индексом, 11.08)
};
\legend{Медиана, P95}
\end{axis}
\end{tikzpicture}
\caption{Время селективной выборки без индекса и с индексом по игроку}
\label{fig:fact-indexes}
\end{figure}

Без индекса выборка выполняется последовательным просмотром: СУБД читает
все 159\,652~строки таблицы фактов и отбрасывает 159\,242~из них,
не удовлетворяющие условию по \texttt{player\_id}. Индекс
\texttt{idx\_gps\_player\_id} переводит запрос на доступ по индексу
(Bitmap Index Scan), при котором читаются только 410~релевантных строк.
Медианное время снижается с 13{,}05 до 2{,}46~мс, P95~--- с 48{,}74
до 11{,}08~мс, что соответствует ускорению примерно в 5{,}3~раза.

\section{Влияние кэширования}

Кэширование проверялось на маршруте \texttt{/stats/leaders} с параметрами
\texttt{metric=per}, \texttt{season\_id=5} и \texttt{limit=20}. В сценарии
Cache Miss перед каждым запросом удалялись ключи лидерборда. В сценарии
Cache Hit кэш предварительно прогревался.

Результаты сравнения приведены в таблице~\ref{tab:cache-effect}.

\begin{table}[H]
\caption{Время ответа API при холодном и тёплом кэше}
\label{tab:cache-effect}
\centering
\small
\begin{tabular}{|p{0.30\textwidth}|p{0.18\textwidth}|p{0.18\textwidth}|p{0.18\textwidth}|}
\hline
\textbf{Сценарий} & \textbf{Медиана (мс)}
  & \textbf{Среднее (мс)} & \textbf{P95 (мс)} \\
\hline
Cache Miss, PostgreSQL & 11{,}68 & 11{,}76 & 15{,}50 \\
\hline
Cache Hit, Redis & 3{,}69 & 4{,}66 & 11{,}57 \\
\hline
\end{tabular}
\end{table}

При Cache Hit медиана ответа снижается с 11{,}68 до 3{,}69~мс,
P95~--- с 15{,}50 до 11{,}57~мс. Ускорение составляет примерно 3{,}2~раза:
API читает готовый JSON из Redis и не выполняет SQL-запрос.

\section{Влияние уровня нагрузки}

Нагрузочное исследование выполнялось для маршрута \texttt{/stats/leaders}
с параметрами \texttt{metric=per}, \texttt{season\_id=5} и \texttt{limit=20}.
Перед каждым раундом выполнялся прогревочный запрос, поэтому замеры отражают
работу API при тёплом Redis-кэше. Обращение без кэша рассмотрено отдельно
в разделе~4.4.

Результаты нагрузочного тестирования приведены в таблице~\ref{tab:load-warm}.

\begin{table}[H]
\caption{Время ответа API при параллельной нагрузке}
\label{tab:load-warm}
\centering
\small
\begin{tabular}{|p{0.24\textwidth}|p{0.16\textwidth}|p{0.16\textwidth}
  |p{0.16\textwidth}|p{0.16\textwidth}|}
\hline
\textbf{Параллельных запросов} & \textbf{Медиана (мс)}
  & \textbf{Среднее (мс)} & \textbf{P95 (мс)} & \textbf{Ошибки (\%)} \\
\hline
1   & 3{,}25   & 3{,}64   & 6{,}62   & 0 \\
\hline
10  & 16{,}43  & 17{,}20  & 32{,}21  & 0 \\
\hline
50  & 71{,}00  & 78{,}61  & 139{,}82 & 0 \\
\hline
100 & 252{,}14 & 326{,}10 & 821{,}08 & 0 \\
\hline
\end{tabular}
\end{table}

Зависимость времени ответа от числа параллельных запросов показана
на рисунке~\ref{fig:load-warm}.

\begin{figure}[H]
\centering
\begin{tikzpicture}
\begin{axis}[
    width=0.88\textwidth,
    height=6.8cm,
    xlabel={Число параллельных запросов},
    ylabel={Время ответа API (мс)},
    xmode=log,
    log basis x=10,
    xtick={1,10,50,100},
    xticklabels={1,10,50,100},
    ymin=0, ymax=860,
    ymajorgrids=true,
    grid style=dashed,
    legend style={
      at={(0.5,1.14)},
      anchor=south,
      legend columns=3,
      font=\small,
      draw=none
    },
    mark size=3pt,
]
\addplot+[mark=*, thick, color=black] coordinates {
    (1, 3.64)
    (10, 17.20)
    (50, 78.61)
    (100, 326.10)
};
\addplot+[mark=square*, thick, color=black, dashed] coordinates {
    (1, 6.62)
    (10, 32.21)
    (50, 139.82)
    (100, 821.08)
};
\addplot+[mark=none, thick, color=black, dotted] coordinates {
    (1, 100) (100, 100)
};
\legend{Среднее время, P95, Требование 100\,мс}
\end{axis}
\end{tikzpicture}
\caption{Время ответа API при росте параллельной нагрузки}
\label{fig:load-warm}
\end{figure}

При нагрузке до 10~запросов P95 (32{,}21~мс) укладывается в требование 100~мс.
При пиковой нагрузке в 100~подключений показатели времени ответа превышают
требования из-за исчерпания пула соединений, однако система сохраняет
отказоустойчивость. Деградация обусловлена насыщением пула соединений
с PostgreSQL и очереди event loop единственного экземпляра Uvicorn.

\section{Анализ результатов и рекомендации}

\textbf{1. Предвычисленная витрина снижает стоимость аналитических выборок.}

На полном объёме прямая агрегация занимает 35{,}58~мс, выборка из витрины~---
1{,}15~мс. Предварительный расчёт сезонных агрегатов уменьшает время
аналитического запроса примерно в 31~раз.

\textbf{2. Индекс по столбцу фильтрации ускоряет селективные выборки.}

Для выборки игрового лога одного игрока индекс \texttt{idx\_gps\_player\_id}
переводит запрос с последовательного просмотра на доступ по индексу
и снижает медианное время примерно в 5{,}3~раза (с 13{,}05 до 2{,}46~мс).

\textbf{3. Redis-кэш снижает задержку одиночных аналитических запросов.}

Переход от Cache Miss к Cache Hit уменьшает медиану ответа с 11{,}68
до 3{,}69~мс и P95 с 15{,}50 до 11{,}57~мс. Кэш следует применять для
повторяющихся рейтингов и дашбордов.

\textbf{4. Порог устойчивой параллельной нагрузки находится между 10 и 50 запросами.}

При 10~одновременных запросах P95 составляет 32{,}21~мс. При 50~запросах
P95 превышает 100~мс (139{,}82~мс), хотя медиана остаётся равной 71{,}00~мс. Для пиковых
нагрузок требуется масштабирование API-сервера: несколько worker-процессов,
балансировщик и общий Redis-кэш.

\section*{Вывод}

Предвычисленная витрина ускоряет аналитические выборки примерно в 31~раз
на полном объёме данных. Индекс \texttt{idx\_gps\_player\_id} ускоряет
селективную выборку по игроку примерно в 5{,}3~раза. Redis-кэш уменьшает
медианную задержку API примерно в 3{,}2~раза. При параллельной нагрузке система не возвращает ошибок,
но требование P95 менее 100~мс выполняется до 10~одновременных запросов
включительно; для большей нагрузки требуется масштабирование API-уровня.
